\section{Savannah's Growth Trajectory: The Decision Window Is Now}
\label{sec:savannah}

\subsection{The Fastest-Growing US Container Port}

The Port of Savannah, operated by the Georgia Ports Authority, processed approximately 5.9 million TEUs in fiscal year 2023---a figure that makes it the fourth-largest container port in the United States by TEU volume and the fastest-growing major port in the country \citep{GPA_annual2023}. Between 2020 and 2024, Savannah's container throughput grew at a compound annual growth rate of approximately 12\%. At that rate, Savannah will overtake the Port of New York/New Jersey (approximately 9.5 million TEUs in 2023) to become the second-largest US container port by approximately 2028.

This growth trajectory is not speculative: it is supported by committed investment. The Georgia Ports Authority has completed a \$2.5 billion infrastructure expansion program (Garden City Terminal expansion, Phase II) that added 3.5 million TEUs of annual capacity to the terminal \citep{GPA_annual2023}. Ocean carriers---including MSC, Maersk, and CMA CGM---have announced or implemented service increases to Savannah as supply chain diversification from West Coast ports accelerates following the 2014--2015 and 2022--2023 ILWU labor disruptions. The port's geographic advantage---24-hour truck gate access, 56-foot channel depth after the \$1.3 billion deepening project, and location at the junction of I-16 and I-95---has made it the primary alternative to LA/Long Beach for East Coast-bound container imports.

\subsection{The Connector Infrastructure Lag}

The connector from the Garden City Terminal to the I-16/I-95 interchange is approximately 6 miles of US-17/SR-25 (Ga.\ Highway 25 Spur), a two-lane undivided state road in several sections, before connecting to I-16. The I-16/I-95 interchange---the junction through which virtually all Savannah port freight must pass---operates at an estimated V/C of 1.31 at peak (Table~\ref{tab:vc}).

The connector's current configuration: US-17 at 4 lanes (widened) from the port gate to the Lynes Parkway interchange, narrowing to 2 undivided lanes through the historic Savannah neighborhoods of West Chatham and Garden City. Peak-hour truck volumes in the 2-lane section reach approximately 2,800 trucks per hour in both directions---more than the design capacity of 2 lanes at HCM standard (approximately 1,600--2,000 vehicles per hour per lane for mixed traffic \citep{HCM6}).

At the current 12\% CAGR, Savannah's port truck traffic will double between 2024 and 2030. The 2-lane section of US-17 operates at V/C 1.31 today; at doubled traffic with no connector expansion, the V/C rises to approximately 2.6---a figure that represents complete breakdown (indefinite queuing, no equilibrium flow state). The standard HPMS V/C calculation for the 6-mile connector as a whole (including the 4-lane sections) yields a 24-hour average V/C of 0.77 (Table~\ref{tab:vc}), which registers as ``adequate capacity''---precisely the measurement error that this paper diagnoses.

\subsection{Investment Lead Times and the Decision Window}

Infrastructure investment in port connector segments follows a predictable timeline. The process from project initiation to construction completion for a major highway connector typically spans:

\begin{itemize}
  \item Environmental review (NEPA): 2--4 years for a major project in an urban/suburban environment
  \item Design and right-of-way acquisition: 2--3 years
  \item Construction: 2--4 years
  \item Total: 6--11 years, with 7--9 years being typical for a project of this scale
\end{itemize}

If Savannah's connector becomes the binding East Coast freight constraint in 2028---as the V/C trajectory implies---and the investment lead time is 7--9 years, then a project initiated in 2026 would complete in 2033--2035---five to seven years after the constraint binds. During those years, port capacity is throttled by connector congestion, not by terminal capacity; the GPA's \$2.5 billion terminal expansion investment is effectively stranded, producing throughput limited by the connector rather than the terminal.

The decision window is now in the sense that an investment decision made today (2026) produces relief in 2033--2035, roughly coinciding with the projected point at which port growth outpaces connector capacity at the projected CAGR. A decision delayed until the constraint visibly binds (2028--2030) produces relief no earlier than 2035--2040---5--10 years of constrained throughput on the East Coast's fastest-growing port.

\subsection{Savannah Investment Scope}

The required investment to relieve the Savannah connector bottleneck involves three components:

\begin{enumerate}
  \item \textbf{US-17/Garden City Connector expansion}: Widen the 2-lane undivided section from Garden City Terminal to the I-16/I-95 interchange from 2 to 6 lanes (matching the widened 4-lane section north of the terminal). Estimated cost: \$620--840 million, including right-of-way acquisition in the constrained West Chatham corridor and two grade-separated interchange replacements \citep{GDOT_2023}.

  \item \textbf{I-16/I-95 interchange reconstruction}: Replace the current diamond interchange with a diverging diamond or turbine interchange capable of handling the projected 2030 peak demand. Estimated cost: \$180--260 million \citep{FHWA_freight2023}.

  \item \textbf{Truck-only bypass lane}: A dedicated truck lane bypassing the signalized section of the connector approach, managed by the GPA as a port access control measure. Estimated cost: \$90--130 million.
\end{enumerate}

Total investment: \$890 million--\$1.23 billion. NPV calculation: Savannah port generates approximately \$2.3 billion per year in direct economic activity and approximately \$19 billion per year in freight value throughput \citep{GPA_annual2023}. A 15\% throughput reduction from connector saturation would cost approximately \$2.85 billion per year. At a 7\% discount rate over 30 years, the NPV of avoiding this loss is approximately \$35 billion---a BCR of 28--40:1 on the \$890 million to \$1.23 billion investment.
